\documentclass[10pt]{article}

\usepackage[utf8]{inputenc}

\usepackage[T1]{fontenc}

\usepackage{amsmath}

\usepackage{amsfonts}

\usepackage{amssymb}

\usepackage[version=4]{mhchem}

\usepackage{stmaryrd}

\usepackage{bbold}

\usepackage{graphicx}

\usepackage[export]{adjustbox}

\graphicspath{ {./images/} }



\begin{document}

ПОСТА МАІІИНА - один вз вариантов Тьюринга машинь.



ПОСТА НОРМАЛЬНАЯ СИСТЕМА, н о р м а л ьн о е и с ч и с л е н и е,- важный частный случай Поста канонической системь. ПОСТА РЕІІЕТКА - решетка по вкЛючению Поста алгебр.



ПОСТА СИСТЕМА ПРОДУКЦИИ, н о р м а л ьная с истема Поста, нормально е и с ч и л л н и е П о с т а,- частный случай Поста канонической системы, когда все правила вывода имеют вид $\frac{G p}{p G}$ и имеется только одно исходное слово (одна аксиома рассматриваемого исчисления). Э. Пост [1] установил әквивалентность П. с. П. и канонич. систем Поста в широком смысле. П. с. П. были использованы Э. Постом и А. А. Марковым (1947) при построении первых примеров ассочиативных исчислений с нераарешимой проблемой распознавания равенства слов (п р о бл е м a T y a).



Лum.: [1] P o st E. L., "Amer. J. Math.", 1943, v. 65, № 2, p. 197-215; [2] М а р к о в А. А., Теория алгорифмов,



М., 1954 (Тр. Матем. ин-та АН СССР, т. 42). С. С. И. Адяле.

ПОСТНИКОВА КВАЛРАТ - коголологическая операция типа $0(1, A, 3, B)$, где $A, B$ - абелевы группы с фиксированным ге те р о м о р ф и 3 м о м $A \rightarrow B$, т. е. таким отображением, что функция



\[

h\left(g_{1}, g_{2}\right)=\eta\left(g_{1}+g_{2}\right)-\eta\left(g_{1}\right)-\eta\left(g_{2}\right)

\]



билинейна и $\eta(-g)=\eta(g)$. Пусть $\xi: F \rightarrow A-$ эпиморфизм, а $\quad F=\bigoplus \mathbb{Z}-$ свободная абелева группа. Для 1-коциклов П. к. определен формулой



\[

e^{1} \longrightarrow \tilde{\eta} \tilde{\xi}\left(e_{0}^{1} \bigcup \delta e_{0}^{1}\right),

\]



где $e_{0}^{1}$ - такая коцепь с коэфффициентами в $F$, что $\tilde{\xi} e_{0}^{1}=e^{1}$. Надстройкой над П. к. является Понтряәина көадрат. Для односвязного $X$ П.к., для к-рого $A=$ $=\pi_{2}(X), B=\pi_{3}(X)$, а $\eta$ определяется композицией с отображением Хопфа $S^{3} \rightarrow S^{2}$, используется при классификации отображений трехмерных полиэдров в $X$. П. к. введен М. М. Постниковым [1].



Лит.: [1] П о с т н и к о в М. М., "Докл. АН СССР», 1949 т. 64, № 4, с. 461-62. А. Ф. Харииладзе IОСТНИКОВА СКСТЕМА, н а т у р а л в н а я система, гомотопическая резольв е т а, П-разложение общего типа,последовательность расслоений



\[

\cdots \stackrel{p_{n+1}}{\longrightarrow} X_{n} \stackrel{p_{n}}{\longrightarrow} X_{n-1} \stackrel{p_{n-1}}{\longrightarrow} \cdots \stackrel{p_{1}}{\longrightarrow} X_{0}=p t

\]



слоями к-рых являются Эйленберга - Маклейна пространства $K\left(\pi_{n}, n\right)$, где $\pi_{n}$ - нек-рая группа (абелева при $n>1$ ). Введена М. М. Постниковым [1]. Пространство $X_{n}$ наз. $n$-м ч л е н о м (или $n-\mathrm{M}$ Э т а ж о м) П. с. $\left\{p_{n}: X_{n} \rightarrow X_{n-1}\right\}$. П. с. $\left\{p_{n}: X_{n} \rightarrow X_{n-1}\right\}$ наз. с х о д я щ е й с я к пространству $X$, если ее обратный предел $\lim \left\{p_{n}: X_{n} \rightarrow X_{n-1}\right\}$ слабо гомотопически эквивалентен пространству $X$. В этом случае пространство $X$ наз. П р е д е л о м П. с. $\left\{p_{n}: X_{n} \rightarrow X_{n-1}\right\}$.



М о р ф и з м о м П. с. $\left\{p_{n}: X_{n} \rightarrow X_{n-1}\right\}$ в П.с. $\left\{q_{n}: Y_{n} \rightarrow Y_{n-1}\right\}$ наз. последовательность непрерывных отображений $f_{n}: X_{n} \rightarrow Y_{n}$, для к-рых диаграмма гомотопически коммутативна. Морфизм $\left\{f_{n}\right\}$ индуцирует отображение $\lim f_{n}: \lim X_{n} \rightarrow \lim Y_{n}$, называемое его п р е пе ло о.



Из определения П. с. следует, что для каждого $n \geqslant 1$ отображение $p_{n}$ является $(n-1)$-эквивалентностью (см. Гомотопическиц̆ тип), в частности $\pi_{i}\left(X_{n-1}\right) \cong$ $\cong \pi_{i}\left(X_{n}\right)$ при $i<n, \pi_{n}\left(X_{n}\right) \cong \pi_{n}$ и $\pi_{i}\left(X_{n}\right)=0$ при $i>n$. Пространства $X$ и $X_{n}$ имеют один и тот же $(n+1)$-тип. В частности, если П. с. конечна, т. е. для нек-рого числа $N$ при всех $n>N$ группа $\pi_{n}$ тривиаль- на, то пространства $X$ и $X_{i}$ гомотопически әквивалентны. В общем случае при $i \leqslant n$ имеютместо изоморфизмы $H_{i}\left(X_{n}\right) \cong H_{i}(X)$ и $\pi_{i}\left(X_{n}\right) \cong \pi_{i}(X)$, т. е. с ростом $n$ группы гомологий и гомотопич. групшы стабилизируются. Для любого кле-



точного разбиения $K \ldots \rightarrow \mathrm{x}_{11} \stackrel{\mathrm{i}_{\cdot} \cdot}{\longrightarrow} \mathrm{X}_{n-1} \rightarrow \cdots \rightarrow \mathrm{X}_{1}$ размерности $\leqslant n \quad$ множества $[K, X]$ и $[K$, $\left.X_{n}\right]$ совпадают. Характеристич. класс $\quad k_{n}=$



$=c\left(p_{n}\right) \in H^{n+1}$$\left(X_{n-1} ;\right.$



\includegraphics[max width=\textwidth, center]{2023_07_08_84f7da0364ad910c7349g-1}

$=c\left(p_{n}\right) \in H^{n+1} \quad\left(X_{n-1} \quad\right.$ Рис. 1,



$\left.\left\{\pi_{n}\right\}\right)$ расслоения $p_{n}: X_{n} \rightarrow X_{n-1}$, т. е. образ при трансгрессии



\[

\tau: H^{n}(K(\pi, n) ; \pi) \longrightarrow H^{n+1}(B ;\{\pi\})

\]



фундаментального класса $\iota_{n} \in H^{n}(K(\pi, n) ; \pi)$, наз. $n$-м $k$-и н в а р и а н т о м (или $n-м$ по с т н и к о в ск и м ф а к т о р о м) П. с. или ее предела $X$. Для любого $n \geqslant 1 n$-й член П. с., а потому и $(n+1)$-тип пространства $X$ полностью определяются группами $\pi_{1}, \ldots$, $\pi_{n}$ и $k$-инвариантами $k_{1}, \ldots, k_{n-1}$. Часто П. с. наз. двойная последовательность



\[

\left\{\pi_{1}, k_{1}, \ldots, \pi_{n}, k_{n}, \ldots\right\} .

\]



Пространство $X$ тогда и только тогда является пределом П. с. $\left\{p_{n}: X_{n} \rightarrow X_{n-1}\right\}$, когда существуют такие $(n-1)$-әквивалентности $\rho_{n}: X \rightarrow X_{n}$, что $\rho_{n-1} \sim$ $\sim \rho_{n} \circ \rho_{n}$ для любого $n \geqslant 1$. Аналогично характеризуются пределы морфизмов П. с.



Существует вариант понятия П. с., иногда оказывающийся более полезным. В этом варианте пространства $X_{n}$ предполагаются клеточными разбиениями, обладающими тем свойством, что $X_{n-1} \subset X_{n}^{n}$ и $X_{n-1}^{n-1}=X_{n}^{n-1}$, а отображения $p_{n}: X_{n} \rightarrow X_{n-1}$ - такими клеточными отображениями (уже не являющимися расслоениями), что, во-первых, $p_{n} \mid X_{n}^{n-1}=\mathrm{id}$ и, во-вторых, гомотопич. слой отображения $p_{n}$ (т. е. слой әтого отображения, превращенного в расслоение) является пространством $K\left(\pi_{n}, n\right)$. Такие П. с. наз. к л е т о ч н ы м и. Пределом клеточной П. с. является клөточное разбиение $X$, для к-рого $X^{n}=X_{n}^{n}$ при любом $n \geqslant 1$. Произвольная П. с. гомотопически эквивалентна клеточной П.с. Основная теорема теории П. с. утверждает (см. [1] [6]), что каждое пространство $X$ является пределом нек-рой однозначно (с точностью до изоморфизма) определенной П. с. $\left\{p_{n}: X_{n} \rightarrow X_{n-1}\right\}$. Эта П. с. наз. системой Постников а прост ра нств а $X$. Имеет место также вариант основной теоремы для отображений: любое отображение $f: X \rightarrow Y$ является пределом нек-рого морфизма $\left\{f_{n}: X_{n} \rightarrow Y_{n}\right\}$ П. с. $\left\{p_{n}: X_{n} \rightarrow X_{n-1}\right\}$ пространства $X$ в П.с. $\left\{q_{n}: Y_{n} \rightarrow Y_{n-1}\right\}$ пространства $Y$. Этот морфизм наз. системо й Постников а о тоб ражения $f$ (др. названия: г о м о то пиче с ка я р езольве н т а о тоб р а ения,П-систем а о б ше го типа отображения, система мураПостни ко в а о то б р а ения). Для постоянного отображения $c: X \rightarrow p t$ линейно связного пространства $X$ его П. с. совпадает с П. с. пространства $X$.



В приложениях большое распространение получили т. н. с т анда р т ны е с ис тем ы П остни ков а (к-рые зачастую наз. просто с и с те ма м и П о с тн и к о в а), представляющие собой П. с., составленные из главных расслоений $p_{n}: X_{n} \rightarrow X_{n-1}$, индуцирующихся из стандартных расслоений Серра $K\left(\pi_{n}, n\right) \rightarrow$ $\rightarrow E K\left(\pi_{n}, n+1\right) \rightarrow K\left(\pi_{n}, n+1\right)$ постниковскими факторами $k_{n} \in H^{n+1}\left(X_{n-1} ; \pi_{n}\right)$, интерпретируемыми в силу представимости групп когомологий как отображения $k_{n}: X_{n-1} \rightarrow K\left(\pi_{n}, n+1\right)$. Стандартными П. с. обладают все пространства, гомотопически простые во всех





\end{document}